\section{Method}
\label{sec:method}

\subsection{Finite-set projection and encoding width}

Let \(\mathcal P\subset\R^d\) be the set of all time-point vectors appearing
in the finite collection of windows to be encoded; \(|\mathcal P|\leq wN\).
For a chosen register width \(n\), draw
\[
    \Pi\in\R^{n\times d},
    \qquad
    \Pi_{ij}\stackrel{\mathrm{i.i.d.}}{\sim}\mathcal N(0,1/n),
\]
and set \(z=\Pi x\). Each coordinate is converted to a bounded rotation angle,
for example \(\theta(z_j)=\pi\tanh(z_j)\), and encoded with a single-qubit
\(R_y(\theta(z_j))\) gate.

The JL guarantee in \cref{prop:jl} concerns only the finite point set
\(\mathcal P\). It permits \(n\) to be chosen from a distortion and failure
probability target instead of setting \(n=d\). Applying the dense projection
still costs \(O(nd)\) arithmetic operations per time point, and neither the
bounded angle map nor the reservoir dynamics is asserted to preserve the JL
geometry.

\subsection{Reset-contractive quantum dynamics}

Sub-reservoir \(r\) uses the input-dependent unitary
\begin{equation}
    V_r(x)
    =
    W_r
    \bigotimes_{j=1}^n R_y\!\left(\theta((\Pi x)_j)\right),
    \label{eq:input-unitary}
\end{equation}
where \(W_r\) is a fixed hardware-topology-aware unitary. One possible block
is an Ising-style layer of edge-wise \(R_{zz}\) gates followed by local
rotations. Its parameters are sampled once and then held fixed.

\begin{figure}[t]
\centering
\resizebox{\linewidth}{!}{
\begin{quantikz}[column sep=0.30cm,row sep=0.22cm]
\lstick{$q_1$} & \gate{R_y(\theta(z_1))} &
  \ctrl{1} & \qw &
  \ctrl{2} &
  \gate{R_z(\vartheta_z^1)} & \gate{R_x(\vartheta_x^1)} & \qw \\
\lstick{$q_2$} & \gate{R_y(\theta(z_2))} &
  \control{} & \ctrl{1} & \qw &
  \gate{R_z(\vartheta_z^2)} & \gate{R_x(\vartheta_x^2)} & \qw \\
\lstick{$q_3$} & \gate{R_y(\theta(z_3))} &
  \qw & \control{} & \control{} &
  \gate{R_z(\vartheta_z^3)} & \gate{R_x(\vartheta_x^3)} & \qw
\end{quantikz}
}
\caption{Example input-dependent unitary for a three-qubit sub-reservoir.
The circuit is an architectural instance, not a claim about a specific
device's native gate set or noise.}
\label{fig:unitary}
\end{figure}

After the unitary, apply the reset channel
\begin{equation}
    T_{r,x}(\rho)
    :=
    \lambda_r V_r(x)\rho V_r(x)^\dagger
    +(1-\lambda_r)\rhoreset,
    \qquad 0\leq\lambda_r<1.
    \label{eq:reset-channel}
\end{equation}
The contraction parameter \(\lambda_r\) controls the loss of initialization
dependence: larger values retain information longer, while smaller values
forget more rapidly. The displayed channel can be dilated with an \(n\)-qubit
reset register and one coin qubit. For the dilation in \cref{fig:reset-dilation},
\(\rhoreset=\lvert+\rangle\!\langle+\rvert^{\otimes n}\). After the controlled
swaps, the ancillary and coin registers are discarded; this partial trace is
essential to obtaining the reduced mixture in \cref{eq:reset-channel}.

\begin{figure}[t]
\centering
\resizebox{0.92\linewidth}{!}{
\begin{quantikz}[column sep=0.34cm,row sep=0.18cm]
\lstick{$q_1$} & \qw & \swap{3} & \qw & \qw & \qw \\
\lstick{$q_2$} & \qw & \qw & \swap{3} & \qw & \qw \\
\lstick{$q_3$} & \qw & \qw & \qw & \swap{3} & \qw \\
\lstick{$a_1$} & \gate{H}\slice{$\lvert+\rangle^{\otimes3}$} & \swap{-3} & \qw & \qw & \qw \\
\lstick{$a_2$} & \gate{H} & \qw & \swap{-3} & \qw & \qw \\
\lstick{$a_3$} & \gate{H} & \qw & \qw & \swap{-3} & \qw \\
\lstick{$c$} & \gate{R_y(\theta_{\lambda})} &
  \ctrl{-6} & \ctrl{-5} & \ctrl{-4} & \qw
\end{quantikz}
}
\caption{A dilation of the reset channel for three reservoir qubits.
Here \(\theta_\lambda=2\arcsin\sqrt{1-\lambda}\). The reduced reservoir
channel is obtained after discarding the coin and ancilla registers.}
\label{fig:reset-dilation}
\end{figure}

For a finite window, the recursion starts from a fixed
\(\rho_{\mathrm{init}}^{(r)}\) and applies \cref{eq:reset-channel} once per
time point. The infinite-history state is the limit obtained by moving this
initialization into the remote past. Existence and uniqueness of that limit
follow from \cref{lem:contraction}.

\subsection{Multiplexing and exact observable features}

The \(R\) sub-reservoirs are separate channel realizations:
\[
    \bigl(\rho^{(1)}_{w,0},\ldots,\rho^{(R)}_{w,0}\bigr).
\]
They can be run in parallel on separate registers or sequentially by reusing
the same physical register. Parallel execution multiplies peak quantum width;
sequential reuse multiplies state preparations and runtime but leaves peak
reservoir width unchanged.

For a family \(\cO\) of local Hermitian observables, the analytic feature map
is the concatenation in \cref{eq:window-feature}. Its dimension is
\(p=R|\cO|\). Exact expectations are used to define the kernel and in the main
generalization theorem. Finite-shot estimates are considered separately in
\cref{prop:shadows,cor:measurement-stability}.

\subsection{Mat\'ern kernel readout}

For \(r\geq0\), smoothness \(\nu>1\), and length scale \(\xi>0\), define
\begin{equation}
    \varphi_{\nu,\xi}(r)
    :=
    \frac{2^{1-\nu}}{\Gamma(\nu)}
    \left(\frac{\sqrt{2\nu}\,r}{\xi}\right)^\nu
    K_\nu\!\left(\frac{\sqrt{2\nu}\,r}{\xi}\right),
    \label{eq:matern-profile}
\end{equation}
where \(K_\nu\) is the modified Bessel function of the second kind. Set
\(\varphi_{\nu,\xi}(0)=1\), and define
\begin{equation}
    \kappa(z,z^{\prime})
    :=
    \varphi_{\nu,\xi}\!\left(\norm{z-z^{\prime}}_2\right).
    \label{eq:matern-kernel}
\end{equation}
The kernel is classical and is evaluated after measurement. Mat\'ern
smoothness is not invoked as evidence of quantum expressivity; the properties
used below are strict positive definiteness and an RKHS section-stability
bound.

Given exact training features \(z_i=\PhiWin(x_i)\), kernel ridge regression
solves
\begin{equation}
    \min_{h\in\cH_\kappa}
    \frac1N\sum_{i=1}^N(h(z_i)-y_i)^2
    +\eta\norm{h}_{\cH_\kappa}^2.
    \label{eq:krr-objective}
\end{equation}
By the representer theorem \citep{ScholkopfSmola2002}, the solution is
\(h(\cdot)=\sum_i\alpha_i\kappa(z_i,\cdot)\), with
\[
    \bm\alpha=(K+N\eta I)^{-1}\bm y.
\]
Naive dense training requires \(O(N^2)\) kernel storage and \(O(N^3)\)
linear-algebra time. These classical costs remain part of the end-to-end
method.

\subsection{Resource summary}

\Cref{tab:resources} distinguishes peak width, feature dimension, and repeated
executions. Big-\(O\) entries describe the stated implementation, not an
optimality or advantage result.

\begin{table}[t]
\caption{Principal resource parameters of \textsc{QuaRK}.}
\label{tab:resources}
\centering
\begin{tabular}{p{0.27\linewidth}p{0.62\linewidth}}
\toprule
Quantity & Accounting \\
\midrule
Encoded reservoir width & \(n\) qubits per sub-reservoir \\
Displayed reset dilation & Up to \(2n+1\) qubits when reset ancillas and coin are simultaneous \\
Reservoir realizations & \(R\); sequential reuse or \(R\)-fold parallel width \\
Exact feature dimension & \(p=R|\cO|\) \\
Projection work & \(O(nd)\) per time point for a dense \(\Pi\) \\
Finite-shot work & \(R\) state preparations per window per snapshot count; see \cref{prop:shadows} \\
Kernel storage/training & \(O(N^2)\) memory and \(O(N^3)\) naive time \\
\bottomrule
\end{tabular}
\end{table}
